\chapter{INTRODUÇÃO}
%\addcontentsline{toc}{chapter}{INTRODUÇÃO}
%%\thispagestyle{plain}

A Introdução constitui-se na apresentação do propósito e do alcance do trabalho, e compõe-se das seguintes fases: o tema (título, assunto escolhido), a definição do problema a ser investigado (a melhor formulação é em forma de pergunta direta, levantamento de questões), a delimitação (afunilamento do assunto, detalhamento, recorte a ser analisado), a hipótese (afirmação antecipada, suposição), a justificativa (contextualização, importância e o porquê da escolha), além de fornecer uma noção geral do conteúdo do restante do trabalho e ainda destacar a linha de pesquisa na qual o relatório está inserido. 
Os objetivos, a justificativa e a metodologia (também denominado como material e métodos) são apresentados separados da introdução.

Apesar de ser o início do trabalho, é a última parte a ser escrita, pois, por apresentar suas linhas gerais, só poderá ser definida quando todas as etapas estiverem delineadas. É a partir da Introdução que as páginas deverão ser numeradas, mas estas mesmas páginas são contadas a partir da Folha de rosto.

%MAIS DICAS
%
%Delimite o tema na introdução de TCC
%Uma boa introdução de TCC deve procurar falar sobre o tema de forma geral e introduzir o leitor ao assunto principal da pesquisa.
%
%Em um texto de cerca de 10 parágrafos, os principais elementos que compõem o TCC devem ser apresentados, tais como o tema, o problema, o objetivo e a metodologia.
%
%No começo da introdução de TCC, apresente o tema que será abordado e explique sua importância.
%
%Nesta parte, seja objetivo e não se aprofunde muito, pois o tema será tratado mais a fundo ao longo do desenvolvimento do TCC.
%
%Um cuidado que deve-se ter na introdução de TCC, é realizar uma boa delimitação do tema, isto é, deixar clara a questão que será investigada mesmo sem se aprofundar no assunto.
%
%É importante fazer essa delimitação para evitar que a introdução de TCC fique muito abrangente.
%
%Não deixando margem para a banca do tcc levantar questionamentos que não foram tratados pela monografia.
%
%Apresente o problema de seu TCC
%Uma boa maneira de delimitar o tema é apresentar a problemática do trabalho.
%
%Você pode fazer isso elaborando um questionamento que serve como propósito para a investigação realizada.
%
%Todo TCC tem o objetivo de responder a uma pergunta ou discutir um tema de forma aprofundada.
%
%Possibilitando assim que se chegue a conclusões ou considerações fundamentadas.
%
%A apresentação do problema nada mais é do que explicar exatamente qual será a pergunta a ser respondida ou tema a ser tratado pelo seu trabalho acadêmico.
%
%Também é a apresentação do problema que deve deixar clara a relevância do tema da sua monografia.
%
%Você deve ressaltar porque o problema que será tratado pelo seu trabalho é importante e deve ser investigado, mostrando assim o valor da sua investigação.
%
%Exponha os objetivos na Introdução do TCC
%Por fim, é essencial apresentar os objetivos do seu TCC na introdução, pois são eles que vão fornecer o direcionamento dado à pesquisa e apresentar  as etapas fundamentais para resolver o problema que está sendo investigado pelo trabalho.
%
%Aqui, você também pode falar sobre algumas hipóteses que foram levantadas, mesmo que de forma superficial, e que serão confirmadas ou refutadas ao final da sua investigação.
%
%Isto ajuda a dar mais foco para o seu TCC e torna mais claro os objetivos do seu trabalho.
%
%Em conjunto com a exposição dos objetivos, devem ser ser apresentados os métodos utilizados para atingi-los.
%
%De forma concisa, escreva sobre a metodologia usada no TCC e identifique o tipo de pesquisa e os instrumentos utilizados para desenvolver seu trabalho.
%
%Para finalizar, não esqueça de usar as informações da introdução no resumo e também no sumário de seu trabalho.
%
%Seguindo os passos acima, você poderá escrever uma boa introdução de TCC para seu trabalho acadêmico de forma simples e estruturada.




\section{Estrutura do trabalho}

Convém que o último parágrafo da introdução seja o detalhamento do conteúdo de cada um dos arquivos subsequentes. Por exemplo: este trabalho está organizado em 6 capítulos, sendo o primeiro a presente introdução. O capítulo \ref{cap1} apresenta aspectos em oftalmologia pertinentes ao contexto para que o leitor possa compreender as retinografias, a fisiologia da retina e sinais oculares. O capítulo \ref{cap2} apresenta fundamentos relevantes em processamento de imagens, descrição de bases de dados contendo imagens de retina, e ainda, fundamentos sobre \textit{roughness index} e transformada de Hough circular. O capítulo \ref{cap3} traz  uma revisão do estado da arte em \tema. Já o capítulo \ref{cap4} apresenta os métodos propostos e o capítulo \ref{cap5} os resultados obtidos. Por fim é apresentada a conclusão. Assuntos de importância menor foram movidos para o apêndice: um capítulo sobre cores e modelos de cores e outro sobre glaucoma e retinopatia diabética.
